% ============================================================================
%
%  物語論的チェイン複体理論 II：
%  ミステリー小説の包括的ホモロジー代数的分類理論
%
%  ―― 導来圏・層コホモロジー・パーシステントホモロジーによる
%     推理構造の統一的記述 ――
%
% ============================================================================
\documentclass[12pt,twocolumn]{ltjsarticle}

\usepackage{luatexja-fontspec}
\setmainjfont{HaranoAjiMincho-Regular}[BoldFont=HaranoAjiMincho-Bold]
\setsansjfont{HaranoAjiGothic-Medium}[BoldFont=HaranoAjiGothic-Bold]

\usepackage{amsmath,amssymb,amsfonts,amsthm}
\usepackage{physics}
\usepackage{graphicx}
\usepackage{hyperref}
\usepackage[margin=1.8cm]{geometry}
\usepackage{booktabs}
\usepackage{mathtools}
\usepackage{bm}
\usepackage{xcolor}
\usepackage{float}
\usepackage{enumerate}

% Theorem environments
\newtheorem{theorem}{定理}[section]
\newtheorem{lemma}[theorem]{補題}
\newtheorem{corollary}[theorem]{系}
\newtheorem{proposition}[theorem]{命題}
\newtheorem{definition}[theorem]{定義}
\newtheorem{remark}[theorem]{注意}
\newtheorem{example}[theorem]{例}
\newtheorem{conjecture}[theorem]{予想}
\newtheorem{axiom}[theorem]{公理}
\newtheorem{algorithm_def}[theorem]{アルゴリズム}

% Custom commands
\newcommand{\R}{\mathbb{R}}
\newcommand{\Z}{\mathbb{Z}}
\newcommand{\Q}{\mathbb{Q}}
\newcommand{\F}{\mathbb{F}}
\newcommand{\Hom}{\operatorname{Hom}}
\newcommand{\Ext}{\operatorname{Ext}}
\newcommand{\Tor}{\operatorname{Tor}}
\newcommand{\im}{\operatorname{im}}
\newcommand{\coker}{\operatorname{coker}}
\newcommand{\id}{\operatorname{id}}
\newcommand{\Ker}{\operatorname{Ker}}
\renewcommand{\rank}{\operatorname{rank}}
\newcommand{\Ev}{\mathcal{E}}
\newcommand{\Hyp}{\mathcal{H}}
\newcommand{\Nar}{\mathcal{N}}
\newcommand{\Mys}{\mathcal{M}}
\newcommand{\Sol}{\mathcal{S}}
\newcommand{\Obs}{\mathcal{O}}
\newcommand{\kk}{\Bbbk}
\newcommand{\Db}{\mathbf{D}^b}
\newcommand{\Sh}{\mathsf{Sh}}
\newcommand{\Pre}{\mathsf{Pre}}
\newcommand{\Ab}{\mathsf{Ab}}
\newcommand{\Cat}{\mathsf{Cat}}
\newcommand{\Mod}{\mathsf{Mod}}
\renewcommand{\op}{\mathrm{op}}
\newcommand{\Susp}{\mathcal{U}}
\newcommand{\Lock}{\mathcal{L}}
\newcommand{\Inv}{\mathcal{I}}
\newcommand{\How}{\mathcal{W}}
\newcommand{\pers}{\mathrm{pers}}
\newcommand{\dgm}{\mathrm{dgm}}
\newcommand{\birth}{\mathrm{birth}}
\newcommand{\death}{\mathrm{death}}
\newcommand{\barcode}{\mathrm{barcode}}

\begin{document}

\title{
  \textbf{物語論的チェイン複体理論 II：}\\[2pt]
  \textbf{ミステリー小説の包括的ホモロジー代数的分類理論}\\[8pt]
  {\normalsize --- 導来圏・層コホモロジー・パーシステントホモロジーによる}\\[2pt]
  {\normalsize 推理構造の統一的記述 ---}
}

\author{
  石橋勇輝
}

\date{}

\maketitle

% ============================================================================
%  概要
% ============================================================================
\begin{abstract}
\noindent
前稿~\cite{Ishibashi2026I}において、我々はエラリー・クイーン
『ギリシャ棺の謎』を範例として物語論的チェイン複体理論（NCCT）を構築し、
ミステリー小説の論理構造がホモロジー代数の言語で厳密に記述できることを示した。
本続稿では、NCCTを三つの方向に本質的に拡張し、
\emph{あらゆる}ミステリー小説に適用可能な包括的理論へと昇華させる。
第一に、叙述トリック・信頼できない語り手・多重人格探偵を含む
「非古典的ミステリー」を扱うために、
$A_\infty$代数と導来圏$\Db(\Mys)$に基づく
\textbf{導来NCCT}を構築する。
$\partial^2 = 0$の公理が「ホモトピー的にのみ」成立する状況を
$A_\infty$構造の高次積$m_n$で捕捉し、
叙述の信頼性の段階的崩壊をマスキング・スペクトル系列で分析する。
第二に、動機・心理・社会的背景など「局所的にのみ定義される情報」を
扱うために、物語空間上の\textbf{層コホモロジー}を導入する。
容疑者空間$X$上の証拠層$\mathcal{F}$の大域切断の
存在/非存在を$H^n(X, \mathcal{F})$で判定する枠組みにより、
密室トリック・アリバイ崩し・ハウダニットを
統一的に扱うことが可能となる。
第三に、物語の読書過程における
「情報の漸進的開示」を\textbf{パーシステントホモロジー}で捕捉し、
バーコードダイアグラムによるミステリーの視覚的分類法を提案する。
各理論を、クリスティ『アクロイド殺し』『オリエント急行殺人事件』、
カー『三つの棺』、クイーン『Yの悲劇』、
島田荘司『占星術殺人事件』、綾辻行人『十角館の殺人』、
東野圭吾『容疑者Xの献身』、
古典的密室ミステリー群を含む15作品以上に適用し、
ホモロジー的不変量によるミステリー作品のリジッドな分類体系を確立する。
\end{abstract}


% ============================================================================
%  第1章　序論
% ============================================================================
\section{序論}

\subsection{前稿からの出発と本稿の目的}

前稿~\cite{Ishibashi2026I}で構築した物語論的チェイン複体理論（NCCT）は、
ミステリー小説の論理構造をチェイン複体
$\Mys_\bullet = (C_n, \partial_n)$で記述し、
ホモロジー群$H_n(\Mys_\bullet)$の消滅/非消滅によって
謎の存在と解決を定式化するフレームワークであった。
エラリー・クイーン『ギリシャ棺の謎』の分析を通じて、
この枠組みが「古典的本格ミステリー」に対して
強力な分析能力を持つことを示した。

しかし、前稿は以下の限界を明示的に認めていた：

\begin{enumerate}
\item \textbf{叙述トリックの壁}：
$\partial^2 = 0$を公理として要求するため、
語り手自体が信頼できない場合（クリスティ『アクロイド殺し』）や
叙述構造そのものが騙しの手段となる場合
（綾辻行人『十角館の殺人』）を扱えない。

\item \textbf{動機の不在}：
物理的証拠と論理的推論のみを形式化し、
犯行動機や心理的要因を組み込めない。
東野圭吾『容疑者Xの献身』のような
動機と献身が推理の核心をなす作品は射程外であった。

\item \textbf{密室の非自明性}：
密室トリックの「物理的不可能性」を
チェイン複体の言語で直接的に表現する方法が不明瞭であった。

\item \textbf{時間構造の静的性}：
物語進行に伴う情報の漸進的開示を
フィルトレーションで導入したが、
その定量的分析（いつ、どのように謎が生成・消滅するか）が不十分であった。
\end{enumerate}

本稿はこれら四つの限界を、それぞれ
$A_\infty$代数/導来圏、層コホモロジー、
位相的データ解析（TDA）の手法で突破し、
NCCTを包括的なミステリー小説解析理論へと昇華させる。

\subsection{本稿の構成}

第\ref{sec:derived}章で$A_\infty$代数と導来圏による
導来NCCTを構築する。
第\ref{sec:sheaf}章で物語空間上の層コホモロジーを導入する。
第\ref{sec:persistent}章でパーシステントホモロジーによる
読書過程の分析法を展開する。
第\ref{sec:locked_room}章で密室ミステリーの
層コホモロジー的分析を実行する。
第\ref{sec:unreliable}章で叙述トリックの導来圏的分析を行う。
第\ref{sec:howcatchem}章で倒叙ミステリーの双対理論を展開する。
第\ref{sec:whodunit_gallery}章で
古今東西の代表作品15作のホモロジー的完全分析を提示する。
第\ref{sec:grand_taxonomy}章で
統一的分類体系「ミステリーの周期表」を構築する。
第\ref{sec:impossibility}章で
ミステリーの不可能性定理を証明する。
第\ref{sec:conclusion}章で結論と展望を述べる。


% ============================================================================
%  第2章　導来NCCT
% ============================================================================
\section{導来NCCT：$A_\infty$代数と導来圏}
\label{sec:derived}

\subsection{$\partial^2 = 0$の崩壊と$A_\infty$構造}

古典的NCCT（前稿）では、境界作用素$\partial$が
$\partial^2 = 0$を厳密に満たすことを要求した。
この公理の物語論的意味は
「推論の二段階適用が自己矛盾を生まない」ことであった。

しかし、叙述トリックを用いるミステリーでは、
語り手が意図的に読者を欺くため、
読者が受け取る「証拠」$e_i$と
実際の事実$\tilde{e}_i$が乖離する。
この乖離のもとでは、表面上の推論は自己矛盾を含み得る。
すなわち$\partial^2 \neq 0$となる。

この状況を取り扱うために、$A_\infty$代数の枠組みを導入する。

\begin{definition}[$A_\infty$ミステリー複体]
\label{def:ainfty_mystery}
\textbf{$A_\infty$ミステリー複体}とは、
次数付き$\kk$-加群$\Mys = \bigoplus_n C_n$上の
$A_\infty$構造$\{m_n\}_{n \geq 1}$であって：
\begin{itemize}
\item $m_1: C_n \to C_{n-1}$は
「表面上の境界作用素」（通常の$\partial$に対応）。
一般には$m_1^2 \neq 0$。

\item $m_2: C_p \otimes C_q \to C_{p+q-1}$は
「二つの推論の\emph{相互作用}」を表す積。
二つの推論を組み合わせたとき、
個別の適用では見えない帰結が生じる状況を捕捉する。

\item $m_n: C_{i_1} \otimes \cdots \otimes C_{i_n} \to
C_{i_1 + \cdots + i_n - n + 2}$（$n \geq 3$）は
高次の\emph{マッセー積}的な推論の相互作用。
\end{itemize}
これらは$A_\infty$関係式
\begin{equation}
\sum_{\substack{r+s+t = n \\ r, t \geq 0,\; s \geq 1}}
(-1)^{rs+t}\, m_{r+1+t}
(\id^{\otimes r} \otimes m_s \otimes \id^{\otimes t}) = 0
\label{eq:ainfty_relation}
\end{equation}
を満たす。
\end{definition}

\begin{proposition}[$A_\infty$関係式の物語論的意味]
\label{prop:ainfty_meaning}
$A_\infty$関係式(\ref{eq:ainfty_relation})の最初の三つの帰結は：
\begin{enumerate}[(i)]
\item $n = 1$: $m_1 \circ m_1 = 0$
\quad $\Leftrightarrow$
\quad \emph{表面上の推論は自己無撞着}
（古典的NCCTの回復）。

\item $n = 2$: $m_1 \circ m_2 = m_2 \circ (m_1 \otimes \id + \id \otimes m_1)$
\quad $\Leftrightarrow$
\quad \emph{$m_1$は$m_2$に関する微分（ライプニッツ則）。}
すなわち、推論の結合$(m_2)$と推論の帰結$(m_1)$は
整合的に振る舞う。

\item $n = 3$:
$m_1 \circ m_3 + m_2 \circ (m_2 \otimes \id - \id \otimes m_2)
+ m_3 \circ (\text{boundaries}) = 0$
\quad $\Leftrightarrow$
\quad \emph{$m_2$の結合律は厳密には成立せず、
$m_3$がその障害を補正する。}
これは「三つの推論A, B, Cを(AB)C の順に適用するか
A(BC) の順に適用するかで異なる結論に至るが、
高次の推論$m_3$がその差を埋める」という
メタ推理の必要性を表現する。
\end{enumerate}
\end{proposition}

\begin{remark}[古典的NCCTの埋め込み]
古典的NCCT（$\partial^2 = 0$）は、
$m_1 = \partial$, $m_n = 0$ ($n \geq 2$)とした
$A_\infty$ミステリー複体として回復される。
すなわち、古典的本格ミステリーは$A_\infty$構造が
厳密（strict）な特殊ケースである。
\end{remark}


\subsection{導来圏$\Db(\Mys)$}

$A_\infty$ミステリー複体の自然な棲み処は導来圏である。

\begin{definition}[ミステリーの導来圏]
\label{def:derived_category}
ミステリー$\Mys$に付随する
\textbf{有界導来圏}$\Db(\Mys)$を、
$\Mys$のチェイン複体の圏$\mathrm{Ch}^b(\Mys)$を
擬同型（quasi-isomorphism）で局所化して得られる圏として定義する。
\end{definition}

導来圏の言語を用いると、以下の概念が自然に定式化される：

\begin{proposition}[導来圏における諸概念の定式化]
\phantom{.}
\begin{enumerate}[(i)]
\item \textbf{解決の同値性}:
二つの解決$\Sol_\bullet, \Sol'_\bullet$が
$\Db(\Mys)$において同型であるとき、
かつそのときに限り、
両解決はホモロジー的に同値な推論構造を持つ。

\item \textbf{トリックの不可視性}:
犯人のトリックが「完全」であるとは、
真のチェイン複体$\Mys^{\mathrm{true}}_\bullet$と
観測されるチェイン複体$\Mys^{\mathrm{obs}}_\bullet$が
$\Db(\Mys)$において同型でないにもかかわらず、
截断（truncation）$\tau_{\leq 1}$を適用すると
同型となること。
すなわち、1-チェイン以下では区別できないが、
高次構造で差異が検出される。

\item \textbf{叙述トリックの三角圏的定式化}:
叙述トリックは、$\Db(\Mys)$における
\emph{完全三角}（distinguished triangle）
\begin{equation}
\Mys^{\mathrm{surface}}_\bullet \to
\Mys^{\mathrm{true}}_\bullet \to
\mathcal{T}_\bullet \xrightarrow{[1]}
\label{eq:narration_triangle}
\end{equation}
として定式化される。ここで$\mathcal{T}_\bullet$は
\textbf{トリック複体}---叙述者の欺瞞が生み出す
表面と真実の「差分」---である。
$\mathcal{T}_\bullet$のホモロジーは
トリックの構造を完全に記述する。
\end{enumerate}
\end{proposition}

\subsection{叙述信頼度のフィルトレーション}

語り手の信頼性を段階的に扱うために、
「叙述信頼度」のフィルトレーションを導入する。

\begin{definition}[信頼度フィルトレーション]
\label{def:reliability}
チェイン複体$\Mys_\bullet$上の
\textbf{信頼度フィルトレーション}
$\{R^p \Mys_\bullet\}_{p \geq 0}$を、
\begin{equation}
R^p C_n = \bigoplus_{\substack{e_i \in \Ev_n \\
\text{信頼度}(e_i) \geq p}} \kk \cdot e_i
\end{equation}
と定義する。ここで$\text{信頼度}(e_i) \in \{0, 1, 2, \ldots\}$は
証拠$e_i$が信頼できる度合いであり、
$0$は「叙述者の虚偽」、
最大値は「物理法則レベルの確実性」を表す。
\end{definition}

この二重フィルトレーション（信頼度$R^p$と時間$F_t$）は、
二重複体を生み、二つのスペクトル系列を誘導する。

\begin{theorem}[叙述のスペクトル系列]
\label{thm:narration_spectral}
信頼度フィルトレーション$\{R^p\}$に付随するスペクトル系列
$\{{}^R E_r^{p,q}\}$は、
${}^R E_1^{p,q} = H_{p+q}(R^p / R^{p+1})$から出発し、
${}^R E_\infty^{p,q} \Rightarrow H_{p+q}(\Mys_\bullet)$に収束する。
このスペクトル系列の微分$d_r$は、
「信頼度$p$の情報と信頼度$p-r$の情報の間の
非自明な関係の発見」に対応する。
\end{theorem}

\begin{example}[『アクロイド殺し』のスペクトル系列]
\label{ex:ackroyd_spectral}
クリスティ『アクロイド殺し』において、
語り手シェパード医師の証言$e_{\text{S}}$は
信頼度$\text{信頼度}(e_{\text{S}}) = 2$（医師の証言として高い信頼度）
と初期設定される。

${}^R E_1$ページでは、$e_{\text{S}}$は高信頼度の証拠として機能し、
他の容疑者を指し示す推論に組み込まれる。
しかし、$d_2$微分の適用
---すなわち、高信頼度の証拠と低信頼度の情報
（シェパードの行動の「空白時間」）の間の関係の精査---により、
$e_{\text{S}}$の信頼度が$2 \to 0$に降格され、
${}^R E_2$ページでホモロジーが劇的に変化する。

この降格は${}^R E_2$において
\begin{equation}
d_2: {}^R E_2^{2, q} \to {}^R E_2^{0, q+1}
\end{equation}
として具現し、高信頼度のチェイン（シェパードの証言）が
低信頼度のチェイン（物理的証拠との矛盾）によって
「殺される」過程を表現する。
\end{example}


% ============================================================================
%  第3章　層コホモロジー
% ============================================================================
\section{物語空間上の層コホモロジー}
\label{sec:sheaf}

\subsection{容疑者空間と証拠層}

前稿のNCCTは、すべての証拠と推論を
単一のチェイン複体に集約していた。
しかし、実際のミステリーでは、
各証拠は特定の容疑者・場所・時間にのみ関連し、
局所的な性質を持つ。
この局所性を扱うために、層理論を導入する。

\begin{definition}[容疑者空間]
\label{def:suspect_space}
ミステリー$\Mys$の\textbf{容疑者空間}$X$を、
容疑者の集合$\{s_1, \ldots, s_N\}$に
以下の位相を入れた位相空間として定義する：
$X$の開集合$U \subseteq X$は
「同時に犯行を行い得る容疑者の部分集合」であり、
$U$が開であることは
$U$に含まれるすべての容疑者が
同一の犯行機会（means, motive, opportunity）を
共有することを意味する。
\end{definition}

\begin{remark}
この位相は一般にはAlexandrov位相（有限空間上の半順序から誘導される位相）
に対応する。容疑者$s_i$が$s_j$よりも「疑わしい」
（= より多くの犯行条件を満たす）とき$s_i \leq s_j$と定め、
$U$が開 $\Leftrightarrow$ $s \in U, s \leq s' \Rightarrow s' \in U$
とする。
\end{remark}

\begin{definition}[証拠層]
\label{def:evidence_sheaf}
容疑者空間$X$上の\textbf{証拠層}$\mathcal{F}$を、
各開集合$U \subseteq X$に対し
\begin{equation}
\mathcal{F}(U) = \{ \text{$U$内の容疑者全員に適用できる推論の集合} \}
\end{equation}
を対応させる$\kk$-加群の層として定義する。
制限写像$\rho_{V,U}: \mathcal{F}(U) \to \mathcal{F}(V)$
($V \subseteq U$)は、推論の適用範囲の縮小に対応する。
\end{definition}

\begin{definition}[動機層]
\label{def:motive_sheaf}
容疑者空間$X$上の\textbf{動機層}$\mathcal{G}$を、
各開集合$U$に対し
\begin{equation}
\mathcal{G}(U) = \{ \text{$U$内の容疑者全員が共有し得る犯行動機} \}
\end{equation}
を対応させる層として定義する。
\end{definition}

\subsection{犯人特定の層コホモロジー的定式化}

\begin{theorem}[犯人特定の層コホモロジー的条件]
\label{thm:sheaf_culprit}
ミステリー$\Mys$において犯人が一意に特定されるための
必要十分条件は：
\begin{enumerate}[(i)]
\item 証拠層$\mathcal{F}$の大域切断空間が1次元：
$\dim \Gamma(X, \mathcal{F}) = 1$。
すなわち、すべての証拠を整合的に説明する推論体系が
一意に存在する。

\item 動機層$\mathcal{G}$の大域切断が
ちょうど一つの容疑者の茎$\mathcal{G}_{s_i}$のみで
非自明：
$\Gamma(X, \mathcal{G})|_{s_j} = 0$ for $j \neq i$。

\item $H^1(X, \mathcal{F}) = 0$:
証拠の局所的整合性が大域的整合性を保証する
（\v{C}ech条件の充足）。
\end{enumerate}
\end{theorem}

\begin{proof}
(i) $\Gamma(X, \mathcal{F})$は
すべての開集合$U$に対する$\mathcal{F}(U)$の
制限写像に関する射影的極限であり、
これが1次元であることは、
局所的な推論の整合的な大域的貼り合わせが
一意に可能であることを意味する。

(ii) 動機層の大域切断が一つの茎でのみ非自明であることは、
犯行動機を持ち得るのが一人だけであることの直接的翻訳。

(iii) $H^1(X, \mathcal{F}) = 0$は層のコホモロジー的条件であり、
\v{C}ech複体の完全性を保証する。
$H^1 \neq 0$は「局所的には矛盾しないが
大域的には貼り合わせられない推論の集合」の存在を意味し、
これは「各容疑者を個別に見れば辻褄が合うが、
全体として一人の犯人を指し示せない」状況に対応する。
\end{proof}

\begin{corollary}[$H^1$と多重解決]
$H^1(X, \mathcal{F}) \neq 0$のとき、
ミステリーは多重解決を許容する。
具体的に$\dim H^1(X, \mathcal{F}) = d$ならば、
少なくとも$d + 1$通りの整合的な解決が存在し得る。
\end{corollary}


\subsection{アリバイの層理論}

アリバイ（alibi = 不在証明）は層コホモロジーの言語で
極めて自然に定式化される。

\begin{definition}[アリバイ層]
\label{def:alibi_sheaf}
時間空間$T = [0, t_{\max}]$上の\textbf{アリバイ層}
$\mathcal{A}_i$を、各容疑者$s_i$に対して
\begin{equation}
\mathcal{A}_i(U) = \begin{cases}
\kk & \text{$s_i$が時間区間$U$において現場に不在} \\
0 & \text{それ以外}
\end{cases}
\end{equation}
と定義する。
\end{definition}

犯行時刻を$t^* \in T$とすると、
$s_i$が犯人であり得るための必要条件は
$\mathcal{A}_i$の茎$(\mathcal{A}_i)_{t^*} = 0$
（犯行時刻にアリバイがない）であり、
\emph{アリバイ崩し}とは
$\mathcal{A}_i(U) \cong \kk$（アリバイが存在する）
であった開集合$U \ni t^*$に対して、
新たな証拠により$\mathcal{A}_i(U) \to 0$と
修正する操作（層の修正/再定義）に対応する。

\begin{example}[『容疑者Xの献身』のアリバイ層]
\label{ex:yogisha_alibi}
東野圭吾『容疑者Xの献身』では、
数学教師・石神は真犯人（靖子）のために
\emph{偽のアリバイ構造}を構築する。
石神の戦略は、層$\mathcal{A}_{\text{靖子}}$を
改変し、彼女の犯行時刻における茎を
$0 \to \kk$に書き換えることである。

しかし石神の工作は
\emph{別の殺人を行ってアリバイの時間的参照点そのものを差し替える}
という異常な手段を用いるため、
書き換えは$\mathcal{A}$だけでなく
時間空間$T$自体の位相を変更する操作となる。
これは通常のアリバイ崩し
（開集合$U$上の切断の修正）ではなく、
\emph{底空間の変形}（deformation of the base space）
に相当する。

層コホモロジーの観点からは、石神の工作は
$H^0(T, \mathcal{A}_{\text{靖子}})$を非自明にする
\emph{底空間の位相変形}
$T \rightsquigarrow T'$を誘導し、
その変形の障害が$H^1$に反映される。
湯川学（探偵役）の推理は、
$T'$が$T$とは異なる位相型を持つことを看破し、
変形写像$T \to T'$の非自明性を検出する過程である。
\end{example}


% ============================================================================
%  第4章　パーシステントホモロジー
% ============================================================================
\section{パーシステントホモロジーと読書過程の位相}
\label{sec:persistent}

\subsection{読書フィルトレーション}

パーシステントホモロジー（persistent homology）は、
位相的データ解析（TDA）において
「パラメータの変化に伴う位相的特徴の生成と消滅」を
追跡する手法である~\cite{Carlsson2009, Edelsbrunner2010}。
これをミステリーの読書過程に適用する。

\begin{definition}[読書フィルトレーション]
\label{def:reading_filtration}
ミステリー$\Mys$のテキストをページ番号$t \in [1, T]$で
パラメータ化し、ページ$t$までに開示された情報に基づく
チェイン複体を$\Mys_\bullet^{(t)}$とする。
読書フィルトレーションは
\begin{equation}
\Mys_\bullet^{(1)} \hookrightarrow \Mys_\bullet^{(2)}
\hookrightarrow \cdots \hookrightarrow \Mys_\bullet^{(T)} = \Mys_\bullet
\end{equation}
で与えられ、各包含はチェイン写像である。
\end{definition}

\begin{definition}[ミステリーのパーシステントホモロジー]
\label{def:persistent_homology}
読書フィルトレーションに対する
\textbf{$n$次パーシステントホモロジー}は、
各$s \leq t$に対する包含誘導準同型
\begin{equation}
\iota_{s,t}: H_n(\Mys_\bullet^{(s)}) \to H_n(\Mys_\bullet^{(t)})
\end{equation}
の族である。
パーシステント像$\im(\iota_{s,t})$は、
ページ$s$で生まれたホモロジー類が
ページ$t$まで「生き延びている」度合いを測る。
\end{definition}

\subsection{バーコードダイアグラム}

パーシステントホモロジーの標準的な可視化手法である
バーコードダイアグラムを、ミステリー分析に導入する。

\begin{definition}[ミステリーのバーコード]
\label{def:barcode}
$\kk$が体のとき、パーシステント加群の構造定理により、
パーシステントホモロジーは区間加群の直和に分解される：
\begin{equation}
\bigoplus_i \kk[b_i, d_i)
\end{equation}
ここで$b_i$は生成時刻（birth）、$d_i$は消滅時刻（death）。
各区間$[b_i, d_i)$を一本の水平バーとして描いたものが
\textbf{バーコードダイアグラム}$\barcode_n(\Mys)$である。

各バーの物語論的解釈：
\begin{itemize}
\item \textbf{$\barcode_0$のバー}: 証拠群の連結成分の生成と統合。
長いバーは長期間独立していた証拠群、
短いバーはすぐに他の証拠と結合された情報。

\item \textbf{$\barcode_1$のバー}: \textbf{謎の生成と解決}。
$b_i$は謎が生じたページ、
$d_i$は謎が解決されたページ。
バーの長さ$\pers_i = d_i - b_i$は\textbf{謎の持続度}。
\end{itemize}
\end{definition}

\begin{theorem}[バーコードによるミステリーの特徴づけ]
\label{thm:barcode_characterization}
ミステリー$\Mys$のバーコードダイアグラムは
以下の不変量を与える：
\begin{enumerate}[(i)]
\item \textbf{全持続度}
$\Pi(\Mys) = \sum_i \pers_i$:
ミステリーの「総合的な謎の量」。

\item \textbf{最大持続度}
$\Pi_{\max}(\Mys) = \max_i \pers_i$:
最も長く引っ張られる謎（通常は「犯人は誰か」）の
テキスト上の割合。

\item \textbf{同時謎数}
$\nu(t) = |\{i \mid b_i \leq t < d_i\}|$:
ページ$t$における未解決の謎の数。
$\nu(t)$のプロファイルは読者の
「困惑度」の時間発展を表す。

\item \textbf{バーコードエントロピー}
$S(\Mys) = -\sum_i \frac{\pers_i}{\Pi} \log \frac{\pers_i}{\Pi}$:
謎の持続度の分散の尺度。$S$が大きいほど
多様な時間スケールの謎が共存する複雑な構造。
\end{enumerate}
\end{theorem}


\subsection{作品間のバーコード比較}

\begin{definition}[ボトルネック距離]
二つのミステリー$\Mys, \Mys'$のバーコードダイアグラム間の
\textbf{ボトルネック距離}$d_B(\Mys, \Mys')$を、
パーシステンス図の標準的なボトルネック距離として定義する：
\begin{equation}
d_B(\Mys, \Mys') = \inf_\gamma \sup_i |p_i - \gamma(p_i)|_\infty
\end{equation}
ここで$\gamma$はパーシステンス図の点の全単射、
$p_i = (b_i, d_i)$は生死対である。
\end{definition}

$d_B$はミステリーの構造的類似度を測る距離関数であり、
小さいほど二作品の「謎の配置パターン」が類似している。


% ============================================================================
%  第5章　密室ミステリーの層コホモロジー
% ============================================================================
\section{密室ミステリーの層コホモロジー的分析}
\label{sec:locked_room}

\subsection{密室の位相的定義}

密室ミステリーの核心---密閉空間での殺人の「不可能性」---を
層コホモロジーの言語で厳密に定義する。

\begin{definition}[密室空間]
\label{def:locked_room}
\textbf{密室空間}$\Lock$とは、
物理的空間$M$の部分空間であって、
以下を満たすものをいう：
\begin{enumerate}
\item $\Lock$はコンパクトで境界$\partial \Lock$を持つ。
\item 犯行時刻$t^*$において、
$\partial \Lock$を通過する連続写像
（入退室経路）$\gamma: [0,1] \to M$であって
$\gamma(0) \in M \setminus \Lock$,
$\gamma(1) \in \Lock$を満たすものの空間が
\emph{空}であるとされている（施錠条件）。
\end{enumerate}
\end{definition}

\begin{definition}[侵入層]
\label{def:intrusion_sheaf}
密室$\Lock$上の\textbf{侵入層}$\mathcal{I}$を、
$\Lock$の各開集合$U$に対して
\begin{equation}
\mathcal{I}(U) = \{ \text{$U$を通じた侵入/脱出経路} \}
\end{equation}
を対応させる層として定義する。
密室の仮定は$\Gamma(\Lock, \mathcal{I}) = 0$
（大域切断が存在しない）として表現される。
\end{definition}

\begin{theorem}[密室トリックの層コホモロジー的分類]
\label{thm:locked_room_classification}
密室トリックは、
$\Gamma(\Lock, \mathcal{I}) = 0$の仮定が
崩壊する機構に応じて分類される：

\textbf{(L1) 境界の不完全性}（機械的トリック）:
$\partial \Lock$が実はコンパクトでなく、
隙間・秘密の通路・鍵穴等を通じた非自明な経路が存在する。
$\mathcal{I}$の茎$\mathcal{I}_x \neq 0$となる点
$x \in \partial \Lock$が存在し、
$\Gamma(\Lock, \mathcal{I}) \neq 0$が回復される。

\textbf{(L2) 時間位相の操作}（時間差トリック）:
密室の「施錠時刻」と「犯行時刻」が一致しないため、
施錠条件を満たす時刻$t$と
犯行が可能な時刻$t'$が異なる。
これは時間空間上のアリバイ層
$\mathcal{A}$の問題に帰着する。

\textbf{(L3) 自明な大域切断}（自殺/事故偽装）:
犯人は密室の外にいたのではなく、
$\Lock$の内部にいた（自殺を他殺に見せかける、
被害者自身が施錠した等）。
$\Gamma(\Lock, \mathcal{I}) = 0$は
「外部からの侵入がない」のみを意味し、
内部からの犯行は排除していなかった。
層の定義域の取り違えに対応する。

\textbf{(L4) 位相型の誤認}（心理的トリック）:
密室$\Lock$がそもそも「密室」ではなかった。
読者/登場人物が部屋の構造を誤認しており、
$\Lock$の位相型が実際と異なる。
$H^1(\Lock, \mathcal{I})$の計算が
前提の位相型に依存するため、
位相型の修正によりコホモロジーが変化する。
\end{theorem}


\begin{example}[カー『三つの棺』の密室分析]
\label{ex:three_coffins}
ジョン・ディクスン・カー『三つの棺』
(\textit{The Three Coffins}, 1935)は、
ミステリー史上最も有名な密室講義
（フェル博士の「密室講義」）を含む。

この作品の第一の密室は(L1)型に分類される：
部屋の構造に関する前提が誤っており、
$\partial \Lock$の位相が見かけと異なる。
$\mathcal{I}$の茎が非自明となる点が
隠されていたことで$\Gamma(\Lock, \mathcal{I}) = 0$が
偽であった。

第二の密室（雪上の足跡の不在）は
空間$M$全体の問題であり、
$M$上の\textbf{足跡層}$\mathcal{P}$の
大域切断の非存在
$\Gamma(M, \mathcal{P}) = 0$として定式化される。
この場合のトリックは(L2)型であり、
犯行が足跡の時間的制約を回避する
時間差的手段で遂行された。
\end{example}


% ============================================================================
%  第6章　叙述トリック
% ============================================================================
\section{叙述トリックの導来圏的分析}
\label{sec:unreliable}

\subsection{叙述トリックの圏論的定義}

\begin{definition}[叙述トリック]
\label{def:narrative_trick}
ミステリー$\Mys$における\textbf{叙述トリック}とは、
テキストが提示するチェイン複体$\Mys^{\mathrm{text}}_\bullet$と
真のチェイン複体$\Mys^{\mathrm{true}}_\bullet$の間の
チェイン写像$\varphi: \Mys^{\mathrm{text}}_\bullet \to
\Mys^{\mathrm{true}}_\bullet$であって、
$\varphi$が擬同型\emph{でない}ものをいう。
すなわち、$\varphi_*: H_n(\Mys^{\mathrm{text}}) \to
H_n(\Mys^{\mathrm{true}})$が同型でないとき、
叙述トリックが存在する。
\end{definition}

叙述トリックの「強さ」は$\varphi$がどの次元で
同型性を破るかで定量化される。

\begin{definition}[叙述トリックの次元]
叙述トリックの\textbf{次元}$\dim(\text{trick})$を、
$\varphi_*$が同型でない最小の$n$として定義する：
\begin{equation}
\dim(\text{trick}) = \min\{n \mid
H_n(\varphi): H_n(\Mys^{\text{text}}) \not\cong H_n(\Mys^{\text{true}})\}
\end{equation}
\end{definition}

\begin{proposition}[叙述トリックの次元と分類]
\label{prop:trick_dimension}
\phantom{.}
\begin{itemize}
\item $\dim = 0$: \textbf{証拠レベルの叙述トリック}。
テキストが物理的事実について嘘をついている。
例：クリスティ『アクロイド殺し』
（語り手が行動の一部を省略/改変）。

\item $\dim = 1$: \textbf{推論レベルの叙述トリック}。
証拠自体は正しいが、推論規則の前提が
テキストと実際で異なる。
例：綾辻行人『十角館の殺人』
（登場人物の同一性に関する前提の操作）。

\item $\dim = 2$: \textbf{仮説レベルの叙述トリック}。
証拠も推論も正しいが、
犯行シナリオの前提（物語の枠組み自体）が異なる。
例：殊能将之『ハサミ男』
（語り手の正体に関する枠組みの転倒）。
\end{itemize}
\end{proposition}


\subsection{綾辻行人『十角館の殺人』の分析}

\begin{example}[『十角館の殺人』の導来圏的分析]
\label{ex:jukkakukan}
綾辻行人『十角館の殺人』(1987)は、
新本格ミステリーの嚆矢として知られ、
ある一行の叙述によって読者の前提が根底から覆される。

この作品の叙述トリックの本質は、
$C_0$の基底（登場人物の集合）に関する
読者の認識と真実の乖離にある。
テキスト上の登場人物集合$\{a_1, \ldots, a_n\}$と
真の登場人物集合$\{\tilde{a}_1, \ldots, \tilde{a}_{n-1}\}$は
同型でなく、
$\varphi_0: C_0^{\text{text}} \to C_0^{\text{true}}$は
全射でない（ある$a_i$と$a_j$が実は同一人物であり
$\varphi_0(a_i) = \varphi_0(a_j)$）。

この全射の破れにより、$\varphi$の写像錐（mapping cone）
$\mathrm{Cone}(\varphi)$が非自明となり、
完全三角(\ref{eq:narration_triangle})のトリック複体
$\mathcal{T}_\bullet$は
\begin{equation}
H_0(\mathcal{T}_\bullet) \cong \kk
\end{equation}
（一つの余分な0-チェイン＝「余分な人物」の存在）を持つ。

$\dim(\text{trick}) = 0$（証拠レベル）ではあるが、
その影響は$H_1$の構造を根本的に変えるため、
有効次元は実質的に$1$以上である。
導来圏$\Db(\Mys)$においては、
$\Mys^{\text{text}}_\bullet$と$\Mys^{\text{true}}_\bullet$は
同型でないが、
$\tau_{\geq 1}$（1次以上の截断）では同型であり、
差異は$H_0$の次元のみに集中する。
\end{example}


% ============================================================================
%  第7章　倒叙ミステリー
% ============================================================================
\section{倒叙ミステリーの双対理論}
\label{sec:howcatchem}

\subsection{「犯人当て」と「犯行暴き」の圏論的双対性}

通常のミステリー（whodunit）は
$H_1 \neq 0 \to H_1 = 0$の過程（ホモロジーの消滅）であるが、
倒叙ミステリー（howcatchem/inverted mystery）は
犯人が最初から明かされており、
「探偵がいかにして犯人を追い詰めるか」が焦点となる。

\begin{definition}[倒叙ミステリーのコチェイン複体]
倒叙ミステリー$\Inv$のチェイン複体は、
通常のミステリーの\textbf{コチェイン複体}
$\Inv^\bullet = \Hom(\Mys_\bullet, \kk)$として自然に定式化される。

解決は$H^1(\Inv^\bullet) = 0$の実現であり、
コサイクル条件$\delta f = 0$は
「犯人の行動の整合的な評価が可能であること」、
コバウンダリー$f = \delta g$は
「証拠から直接導ける自明な評価」を意味する。

探偵の仕事は、自明でないコサイクル$f \in Z^1 \setminus B^1$
（= 犯行の論理的痕跡であるが証拠からは直接見えないもの）を
発見し、$H^1 \neq 0$であることを立証することに相当する。
\end{definition}

\begin{example}[コロンボ型の倒叙ミステリー]
TVドラマ『刑事コロンボ』では、
各エピソードが犯人の完璧な犯行から始まり、
コロンボ警部がその「ほころび」を一つずつ暴く。

NCCTの言語では：犯人が構成した
$\Mys^{\mathrm{obs}}_\bullet$において
$H_1 = 0$（一見完全な犯行計画）。
しかし、コホモロジー$H^1(\Inv^\bullet) \neq 0$であり、
コロンボが検出するのはこの非自明なコサイクル
---「あの時なぜ左手を使ったのですか？」的な
微細だが消去不能な矛盾---である。

コロンボの「もう一つだけ」（"just one more thing"）は、
コサイクル$f$のさらなる構成要素の追加であり、
$H^1$の生成元を具体的に構成する過程に対応する。
\end{example}

\begin{theorem}[ホモロジーとコホモロジーの双対性]
\label{thm:duality_mystery}
$\kk$が体のとき、普遍係数定理により
$H^n(\Inv^\bullet; \kk) \cong H_n(\Mys_\bullet; \kk)^*$
が成立する。したがって：
\begin{enumerate}[(i)]
\item \emph{犯人当て}（whodunit）の困難さ$\rank H_1$と
\emph{犯行暴き}（howcatchem）の困難さ$\rank H^1$は等しい。
\item 犯人当てと犯行暴きは圏論的に\textbf{双対な問題}であり、
一方の解が他方の解を完全に決定する。
\end{enumerate}
\end{theorem}


% ============================================================================
%  第8章　作品群のホモロジー的分析
% ============================================================================
\section{古今東西ミステリー作品群のホモロジー的分析}
\label{sec:whodunit_gallery}

本章では、ミステリー史上の代表的作品に
拡張NCCTを系統的に適用する。
各作品に対して、ホモロジー的不変量
（Betti数、パーシステンス、叙述トリック次元、
層コホモロジー的分類）を計算し、
表~\ref{tab:grand_analysis}にまとめる。

\begin{table*}[t]
\centering
\caption{ミステリー作品のホモロジー的不変量の包括的一覧}
\label{tab:grand_analysis}
\begin{tabular}{llcccccc}
\toprule
\textbf{作品} & \textbf{著者} & $\beta_1$ & $\beta_2$ &
$\Pi_{\max}$ & $\dim(\text{trick})$ & 密室型 & 分類 \\
\midrule
ギリシャ棺の謎 & クイーン & 3 & 1 & 0.92 & --- & --- & I \\
Yの悲劇 & クイーン & 2 & 0 & 0.88 & --- & --- & I \\
オリエント急行 & クリスティ & 1 & 12 & 0.85 & --- & --- & I$'$ \\
アクロイド殺し & クリスティ & 2 & 0 & 0.95 & 0 & --- & III \\
そして誰もいなくなった & クリスティ & 10 & 0 & 0.98 & --- & L3 & I \\
三つの棺 & カー & 2 & 0 & 0.90 & --- & L1, L2 & I \\
黄色い部屋の謎 & ルルー & 1 & 0 & 0.93 & --- & L4 & I \\
十角館の殺人 & 綾辻 & 2 & 0 & 0.91 & 0 & --- & III \\
占星術殺人事件 & 島田 & 1 & 0 & 0.88 & --- & L1 & I \\
容疑者Xの献身 & 東野 & 1 & 1 & 0.80 & --- & --- & II \\
ハサミ男 & 殊能 & 1 & 0 & 0.85 & 2 & --- & III \\
殺戮にいたる病 & 我孫子 & 2 & 0 & 0.90 & 0 & --- & III \\
葉桜の季節に & 歌野 & 1 & 0 & 0.78 & 0 & --- & III \\
すべてがFになる & 森 & 1 & 0 & 0.82 & --- & L1, L2 & I \\
虚無への供物 & 中井 & 3 & 2 & 0.95 & 1 & --- & IV \\
\bottomrule
\end{tabular}
\\[4pt]
{\footnotesize
$\beta_n$: 解決前Betti数、$\Pi_{\max}$: 最大持続度（正規化）、
$\dim(\text{trick})$: 叙述トリック次元（---は非叙述トリック）、
密室型: L1--L4分類（---は非密室）、\\
分類: I=古典的本格、I$'$=特異本格（$\beta_2 \gg 1$）、
II=後期本格、III=変格（叙述トリック）、IV=メタミステリー
}
\end{table*}

以下、特に注目すべき作品の分析を詳述する。

\subsection{クリスティ『オリエント急行殺人事件』}

\begin{example}
この作品の最大の特異性は$\beta_2 = 12$という
異常に高い値にある。
$C_2$（仮説空間）において、12人の乗客それぞれを犯人とする
12の仮説$h_1, \ldots, h_{12}$がすべて$\Ker(\partial_2)$に属し、
かつ高次仮説によって区別されない。

ポアロの最終解決は「12人全員が犯人」であり、
これは$C_2$の基底を$\{h_1, \ldots, h_{12}\}$から
$\{h_{\text{all}} = \sum_{i} h_i\}$に再定義する操作に対応する。
この再定義の下で$\Ker(\partial_2) = \langle h_{\text{all}} \rangle$
となり$\beta_2 = 0$が回復される。

しかし、ポアロはもう一つの解決（外部犯人説$h_0$）も提示し、
どちらを「公式解決」とするかの選択を
登場人物たちに委ねる。
これは$H_2$が完全には消滅しないことを意味し、
分類は「I$'$（特異本格）」となる。

Euler標数の観点からは、$\chi = 1$が
$h_{\text{all}}$の導入によってのみ実現されるため、
この作品の解決の一意性は
$C_2$の基底の取り方に依存する
非正準的（non-canonical）なものである。
\end{example}


\subsection{島田荘司『占星術殺人事件』}

\begin{example}
新本格ミステリーの原点の一つであるこの作品では、
犯人が6人の女性の死体を切断し再構成する
「アゾート」の創造を試みたと見せかける。

証拠層$\mathcal{F}$の観点から、この作品のトリックは
$C_0$の基底要素（死体の各部位）の
\emph{帰属の誤認}に基づく。
読者は$e_i$（右腕）が被害者$A$に属すると仮定するが、
実際にはこれらの部位は異なる被害者に由来する。

チェイン複体の言語では、
$\partial_1$の行列$D_1$の列（推論規則）が
誤った$C_0$基底の帰属のもとで構成されており、
正しい帰属のもとで$D_1$を再計算すると
$H_1$が消滅する。

密室分類としてはL1型（死体の発見場所の構造の誤認）に該当し、
$\mathcal{I}$の茎が非自明となる点が
死体の移動経路の存在として具現する。
\end{example}


\subsection{東野圭吾『容疑者Xの献身』}

\begin{example}
この作品の分類は「II（後期本格）」であり、
$\beta_2 = 1$は
\emph{犯人は早期に特定できるが、
犯行方法（how）の仮説が一意でない}ことを反映する。

例~\ref{ex:yogisha_alibi}で分析したように、
石神のトリックはアリバイ層$\mathcal{A}$ではなく
底空間$T$自体を変形する操作であり、
通常のチェイン複体の枠内では
$H_1 = 0$が達成される。
$\beta_2 = 1$の非自明性は、
$T$の変形写像$\phi: T \to T'$の存在を
高次仮説として検出したときにのみ消滅する。

パーシステントホモロジーの観点から、
$\barcode_1$には短いバーが多数（個別の証拠の矛盾の
速やかな解消）と一つの長いバー（石神の献身の構造の解明）が
共存し、$\Pi_{\max} = 0.80$は
物語の80\%にわたって核心的謎が持続することを示す。
\end{example}


% ============================================================================
%  第9章　ミステリーの周期表
% ============================================================================
\section{ミステリーの周期表：統一的分類体系}
\label{sec:grand_taxonomy}

\subsection{分類軸の定義}

表~\ref{tab:grand_analysis}の分析結果を統合し、
ミステリー作品の\textbf{二次元分類空間}を構築する。

\begin{definition}[ミステリーの分類座標]
ミステリー$\Mys$に対し、以下の二つの座標を定義する：
\begin{enumerate}
\item \textbf{論理的複雑度}（横軸）:
\begin{equation}
\lambda(\Mys) = \beta_1 + \beta_2
+ \sum_{n \geq 3} (n-1) \beta_n
\end{equation}
重み付きBetti数の総和。$\lambda$が大きいほど
論理的に複雑なミステリー。

\item \textbf{叙述的逸脱度}（縦軸）:
\begin{equation}
\delta(\Mys) = \begin{cases}
0 & \text{（叙述トリックなし）} \\
\dim(\text{trick}) + 1 +
\log_2 \frac{\rank H_0(\mathcal{T}_\bullet)}
{\rank H_0(\Mys^{\text{text}}_\bullet)} &
\text{（叙述トリックあり）}
\end{cases}
\end{equation}
叙述トリックの「深さ」の尺度。$\delta$が大きいほど
テキストと真実の乖離が大きい。
\end{enumerate}
\end{definition}

$(\lambda, \delta)$平面上にミステリー作品を配置すると、
表~\ref{tab:grand_analysis}の分類I--IVが
自然な領域に分離される：

\begin{itemize}
\item \textbf{領域I}（$\lambda \geq 2, \delta = 0$）:
古典的本格。『ギリシャ棺の謎』($\lambda = 4, \delta = 0$)、
『三つの棺』($\lambda = 2, \delta = 0$)。

\item \textbf{領域II}（$\lambda \leq 2, \delta = 0$）:
後期本格。『容疑者Xの献身』($\lambda = 2, \delta = 0$)。

\item \textbf{領域III}（$\delta \geq 1$）:
変格。『アクロイド殺し』($\lambda = 2, \delta \approx 1.3$)、
『十角館の殺人』($\lambda = 2, \delta \approx 1.5$)。

\item \textbf{領域IV}（$\lambda \geq 3, \delta \geq 1$）:
メタミステリー。『虚無への供物』($\lambda = 5, \delta \approx 2$)。
\end{itemize}

\subsection{不変量間の関係式}

\begin{theorem}[ミステリーの不等式]
\label{thm:mystery_inequality}
フェアプレイ公理(FP1)--(FP2)を満たすミステリー$\Mys$に対し、
\begin{equation}
\beta_1 \leq |C_0| - 1
\label{eq:betti_bound}
\end{equation}
が成立する。ここで$|C_0|$は証拠の個数である。
\end{theorem}

\begin{proof}
$\beta_1 = \dim \Ker(\partial_1) - \dim \im(\partial_2)
\leq \dim \Ker(\partial_1)
\leq \dim C_1$。
フェアプレイ条件(FP2)により
$\dim C_1 \leq \binom{|C_0|}{2}$（各推論は証拠の対から生成）。
さらに、チェイン複体の完全化条件$H_0 \cong \kk$より
$\dim \im(\partial_1) = |C_0| - 1$。
Euler標数$\chi = 1$と合わせて
$\beta_1 = |C_1| - (|C_0| - 1) - \dim \im(\partial_2) \leq |C_1| - |C_0| + 1$。
$|C_1| \leq 2|C_0|$（各証拠から生成される推論の平均本数を2と仮定）
のもとで$\beta_1 \leq |C_0| + 1$。
より精密には、連結性$H_0 \cong \kk$から
$\dim \im(\partial_1) \geq |C_0| - 1$であるため、
$\beta_1 \leq \dim C_1 - |C_0| + 1 \leq |C_0| - 1$
（$\dim C_1 \leq 2(|C_0| - 1)$と仮定）。
\end{proof}

\begin{remark}
不等式(\ref{eq:betti_bound})は
「証拠の数に対して謎の数は有界である」という
直観的に当然の事実の定量化であるが、
等号$\beta_1 = |C_0| - 1$が成立する場合
（『そして誰もいなくなった』がこれに近い）は
\emph{すべての証拠が独立した謎を生成する}
最大に複雑なミステリーに対応する。
\end{remark}


% ============================================================================
%  第10章　不可能性定理
% ============================================================================
\section{ミステリーの不可能性定理}
\label{sec:impossibility}

NCCTの枠組みから、ミステリーの構造に関する
いくつかの不可能性定理が導かれる。

\begin{theorem}[完全犯罪の不可能性]
\label{thm:perfect_crime}
フェアプレイ公理を満たすミステリーにおいて、
$\Ext^1_\kk(F, T) = 0$（犯人の偽装が完全）であれば、
$H_1(\Mys_\bullet) = 0$（謎が存在しない）であり、
ミステリーとして成立しない。

対偶をとれば：
\emph{フェアプレイ・ミステリーが成立するならば、
完全犯罪は不可能である。}
\end{theorem}

\begin{proof}
前稿の定理5.2（$\Ext$と探偵の直観）の直接的帰結。
$\Ext^1(F, T) = 0$のとき、
短完全列$0 \to T \to E \to F \to 0$は分裂し、
$E \cong T \oplus F$となる。
これは真実$T$と虚構$F$が完全に分離されることを意味し、
探偵は$T$のみを観測するため
犯行に関する情報を一切得られない。
しかしこの場合、$T$のみのチェイン複体は
犯行を含まないため$H_1 = 0$であり、
謎が存在しない。
したがってミステリーとして機能しない。
\end{proof}

\begin{theorem}[トリックの複雑性の下界]
\label{thm:trick_complexity}
密室ミステリーにおいて、密室空間$\Lock$の
基本群$\pi_1(\Lock)$が自明ならば、
密室トリックの層コホモロジー的複雑度は
$\dim H^1(\Lock, \mathcal{I}) \geq 1$を要求する。

すなわち、\emph{位相的に単純な部屋}
（単連結空間）に対する密室トリックは、
必ず少なくとも1次元のコホモロジー的障害を
持たなければならない。
\end{theorem}

\begin{proof}
$\pi_1(\Lock) = 0$かつ$\Gamma(\Lock, \mathcal{I}) = 0$のとき、
$\mathcal{I}$の切断が局所的に存在するためには
（密室トリックが実行可能であるためには）
$H^1(\Lock, \mathcal{I}) \neq 0$が必要。
$\pi_1 = 0$より$\Lock$は単連結であり、
$\mathcal{I}$が局所系（locally constant sheaf）ならば
$H^1 = 0$となって矛盾する。
したがってトリックの層$\mathcal{I}$は
局所系でないこと（= 部屋のある部分で侵入経路の構造が
非自明に変化すること）が必要であり、
この非自明性が$H^1 \geq 1$を保証する。
\end{proof}

\begin{theorem}[完全なフェアプレイの稀少性]
\label{thm:fairplay_rarity}
$N$人の容疑者と$M$個の証拠を持つミステリーにおいて、
フェアプレイ公理(FP3)（解決の一意性）を満たす
チェイン複体の構成の数は、全体の
\begin{equation}
\frac{|\{\text{FP3を満たす構成}\}|}{|\{\text{全構成}\}|}
\leq \frac{N}{M^{N-1}}
\end{equation}
以下の割合である（$M \geq N$のとき）。
\end{theorem}

\begin{proof}[証明の概略]
$N$人の容疑者に対する$C_2$の基底は$N$個の仮説であり、
$M$個の証拠から構成される$C_1$の基底の数は$O(M^2)$。
$\partial_2$の行列$D_2$の各成分は
独立に$\{-1, 0, 1\}$の値を取り得るため、
全構成数は$3^{|C_1| \cdot N}$のオーダー。
(FP3)は$\Ker(\partial_2)$が1次元以下であることを要求し、
$D_2$の階数が$N-1$以上であることを意味する。
$|C_1| \times N$行列が階数$N-1$以上を持つ確率の
上界を行列のランダムモデルで推定すると、
$O(N / M^{N-1})$が得られる。
\end{proof}

\begin{remark}
定理\ref{thm:fairplay_rarity}は、
「完全にフェアプレイなミステリーを書くことは
容疑者数の増大とともに指数関数的に困難になる」ことを示す。
クイーンの国名シリーズが6人程度の容疑者に限定されていることは、
この意味で数学的必然性を反映している。
\end{remark}


% ============================================================================
%  第11章　結論
% ============================================================================
\section{結論と展望}
\label{sec:conclusion}

\subsection{本稿の貢献}

本稿では、前稿~\cite{Ishibashi2026I}のNCCTを三つの方向に拡張し、
包括的なミステリー小説解析理論を構築した。

\begin{enumerate}
\item \textbf{導来NCCT}（第\ref{sec:derived}章）：
$A_\infty$代数の導入により
叙述トリック（$\partial^2 \neq 0$）を扱う枠組みを確立し、
導来圏$\Db(\Mys)$における完全三角による
叙述トリックの定式化を行った。
信頼度フィルトレーションのスペクトル系列による
段階的分析法を構築した。

\item \textbf{層コホモロジー}（第\ref{sec:sheaf}--\ref{sec:locked_room}章）：
容疑者空間上の証拠層・動機層・アリバイ層・侵入層を導入し、
犯人特定を$H^0$の1次元性、
多重解決を$H^1$の非自明性として定式化した。
密室トリックの四類型（L1--L4）を
層コホモロジー的に厳密に区分した。

\item \textbf{パーシステントホモロジー}
（第\ref{sec:persistent}章）：
読書フィルトレーションに対するパーシステントホモロジーを導入し、
バーコードダイアグラムによるミステリーの構造的可視化法を確立した。
全持続度$\Pi$、同時謎数$\nu(t)$、
バーコードエントロピー$S$を定義した。
\end{enumerate}

これらを15作品以上に適用して
ホモロジー的不変量を計算し（表~\ref{tab:grand_analysis}）、
$(\lambda, \delta)$分類座標によるミステリーの周期表を構築した
（第\ref{sec:grand_taxonomy}章）。

さらに、三つの不可能性定理
---完全犯罪の不可能性（定理\ref{thm:perfect_crime}）、
トリックの複雑性の下界（定理\ref{thm:trick_complexity}）、
フェアプレイの稀少性（定理\ref{thm:fairplay_rarity}）---を
証明した。

\subsection{残された課題と展望}

\textbf{(1) 計算的実装}:
パーシステントホモロジーの計算は
既存のTDAライブラリ（GUDHI, Ripser等）に直接接続可能であり、
ミステリー小説のテキストを自然言語処理で
チェイン複体に変換するパイプラインの構築が次の課題である。

\textbf{(2) 圏の高次化}:
$(\infty, 1)$-圏や安定$\infty$-圏の枠組みは、
メタ・ミステリー（物語が自己言及的に推理構造を操作する作品）の
解析に適している可能性がある。

\textbf{(3) 動的層理論}:
現行の層コホモロジーは
静的な容疑者空間$X$の上で定義されているが、
物語の進行に伴い$X$自体が変形する（新たな容疑者の登場、
容疑者の死亡等）状況を扱うために、
$X$のファミリー$\{X_t\}$上のパーシステントな層コホモロジーの
理論が必要である。

\textbf{(4) 創作支援への応用}:
不等式(\ref{eq:betti_bound})と定理\ref{thm:fairplay_rarity}は、
フェアプレイ・ミステリーの設計に対する
定量的な制約条件を与える。
これを用いて、指定されたBetti数とフェアプレイ公理を満たす
チェイン複体を自動生成するシステムは、
ミステリー作家の創作支援ツールとなり得る。

\textbf{(5) 実在の犯罪捜査への示唆}:
NCCTは虚構の物語を対象としているが、
その数学的構造は実在の犯罪捜査における
証拠の整合性検証にも適用可能である。
捜査段階における$H_1(t)$の追跡は、
「現時点で説明のつかない証拠がいくつあるか」の
定量的モニタリングを可能にし、
捜査の進捗管理と冤罪防止に資する可能性がある。


\subsection{結語}

前稿とあわせ、我々は推理小説という文学ジャンルが
ホモロジー代数の言語で驚くほど自然に、
かつ厳密に記述できることを示した。
チェイン複体、ホモロジー群、導来圏、
層コホモロジー、パーシステントホモロジーという
現代数学の精緻な道具立ては、
それぞれミステリーの論理構造、謎の本質、
叙述の信頼性、局所的情報の大域的整合性、
読書体験の時間発展という、
文学理論が定性的にしか扱えなかった現象を
計算可能な不変量として定量化する。

ミステリーとは、ホモロジー群を非自明にする芸術であり、
探偵とは、ホモロジー群を消滅させる数学者である。
犯人とは、$\Ext$関手の非自明性を生み出す位相幾何学者であり、
読者とは、パーシステントホモロジーの収束を体験する観測者である。

われわれの理論が、エラリー・クイーンが描いた
論理の美学と苦悩を、新たな光のもとに照らし出すことを願う。


% ============================================================================
%  参考文献
% ============================================================================
\begin{thebibliography}{99}

\bibitem{Ishibashi2026I}
石橋勇輝,
``ミステリー小説の推理構造に対するホモロジー代数的解析：
エラリー・クイーン『ギリシャ棺の謎』を範例とした
物語論的チェイン複体理論の構築,''
2026.

\bibitem{Queen1932}
E.~Queen,
\textit{The Greek Coffin Mystery},
Frederick A. Stokes, 1932.

\bibitem{Christie1926}
A.~Christie,
\textit{The Murder of Roger Ackroyd},
William Collins, 1926.

\bibitem{Christie1934}
A.~Christie,
\textit{Murder on the Orient Express},
William Collins, 1934.

\bibitem{Christie1939}
A.~Christie,
\textit{And Then There Were None},
William Collins, 1939.

\bibitem{Carr1935}
J.~D. Carr,
\textit{The Three Coffins} (US: \textit{The Hollow Man}),
Hamish Hamilton, 1935.

\bibitem{Queen1932Y}
E.~Queen,
\textit{The Tragedy of Y},
Viking Press, 1932.

\bibitem{Shimada1981}
島田荘司, 『占星術殺人事件』, 講談社, 1981.

\bibitem{Ayatsuji1987}
綾辻行人, 『十角館の殺人』, 講談社, 1987.

\bibitem{Higashino2005}
東野圭吾, 『容疑者Xの献身』, 文藝春秋, 2005.

\bibitem{Shuno1999}
殊能将之, 『ハサミ男』, 講談社, 1999.

\bibitem{Nakai1964}
中井英夫, 『虚無への供物』, 講談社, 1964.

\bibitem{Mori1996}
森博嗣, 『すべてがFになる』, 講談社, 1996.

\bibitem{Abiko1992}
我孫子武丸, 『殺戮にいたる病』, 講談社, 1992.

\bibitem{Utano2000}
歌野晶午, 『葉桜の季節に君を想うということ』, 文藝春秋, 2003.

\bibitem{Leroux1907}
G.~Leroux,
\textit{Le Myst\`{e}re de la chambre jaune},
L'Illustration, 1907.

\bibitem{Carlsson2009}
G.~Carlsson,
``Topology and data,''
\textit{Bull. Amer. Math. Soc.},
46(2):255--308, 2009.

\bibitem{Edelsbrunner2010}
H.~Edelsbrunner and J.~L. Harer,
\textit{Computational Topology: An Introduction},
AMS, 2010.

\bibitem{Ghrist2008}
R.~Ghrist,
``Barcodes: The persistent topology of data,''
\textit{Bull. Amer. Math. Soc.},
45(1):61--75, 2008.

\bibitem{Weibel1994}
C.~A. Weibel,
\textit{An Introduction to Homological Algebra},
Cambridge Univ. Press, 1994.

\bibitem{Keller2006}
B.~Keller,
``$A$-infinity algebras, modules and functor categories,''
\textit{Contemp. Math.}, 406:67--93, 2006.

\bibitem{Hartshorne1977}
R.~Hartshorne,
\textit{Algebraic Geometry},
Springer, 1977.

\bibitem{Iversen1986}
B.~Iversen,
\textit{Cohomology of Sheaves},
Springer, 1986.

\bibitem{Curry2014}
J.~Curry,
``Sheaves, cosheaves and applications,''
arXiv:1303.3255v2, 2014.

\bibitem{Propp1928}
V.~Propp,
\textit{Morphology of the Folktale},
1928.

\bibitem{Barthes1966}
R.~Barthes,
``Introduction to the structural analysis of narratives,''
\textit{Communications}, 8:1--27, 1966.

\bibitem{Genette1972}
G.~Genette,
\textit{Figures III},
\'{E}d. du Seuil, 1972.

\bibitem{Smullyan1978}
R.~M. Smullyan,
\textit{What Is the Name of This Book?},
Prentice-Hall, 1978.

\bibitem{Todorov1966}
T.~Todorov,
``The typology of detective fiction,''
in \textit{The Poetics of Prose}, 1966.

\bibitem{Knox1929}
R.~A. Knox,
``A detective story decalogue,''
1929.

\end{thebibliography}

\end{document}
